% ch04_u05 —— 原书 p.162–p.171（4.3.17–4.4.16）
% 承接 p.161 末句：……其中 $\tilde{SO}_n(K)$ 是中心扩张
\SourcePage{162}

(196)（对应于 $\omega_{O}^{1/2}$）。$\operatorname{Aut}_{2}O$ 在 $\operatorname{Pf}_{\mathcal{L}}^{\wedge}$ 上的作用是清楚的：$\operatorname{Aut}_{2}O$ 作用在概形 $M_{2}^{\wedge}$ 上，而 $M_{2}^{\wedge}$ 看作 $\operatorname{Bun}_{SO_{n}}$ 上的概形。另一方面，$\tilde{SO}_{n}(K)$ 作用在 $\operatorname{Pf}_{\mathcal{L},\hat x}$ 上，后者定义为 $\operatorname{Pf}_{\mathcal{L}}^{\wedge}$ 在 $M_{2}^{\wedge}$ 于 $\hat x\in X_{2}^{\wedge}$ 处纤维上的限制。事实上，该纤维等于 $\operatorname{Bun}_{SO_{n},\underline{x}}$，其中 $x$ 是 $\hat x$ 在 $X$ 中的像；由 4.3.7，中心扩张 $\tilde{SO}_{n}(K_{x})_{\mathcal{L}}$ 作用在 $\operatorname{Pf}_{\mathcal{L}}$ 到 $\operatorname{Bun}_{SO_{n},\underline{x}}$ 的拉回上。该扩张只依赖 $\mathcal{L}_{x}$，即 $\mathcal{L}$ 到 $\operatorname{Spec}O_{x}$ 的拉回。因 $\hat x$ 给出 $(O,\omega_{O}^{1/2})$ 与 $(O_{x},H^{0}(\operatorname{Spec}O_{x},\mathcal{L}_{x}))$ 之间的同构，我们得到同构 $\tilde{SO}_{n}(K_{x})_{\mathcal{L}}\xrightarrow{\ \sim\ }\tilde{SO}_{n}(K)$，从而得到所要的 $\tilde{SO}_{n}(K)$ 作用。

\bdp{4.3.17} \emph{命题.}
\begin{enumerate}[label=(\roman*),leftmargin=2.2em]
\item $\tilde{SO}_{n}(K)$ 在 $\operatorname{Pf}_{\mathcal{L},\hat x}$（$\hat x\in X_{2}^{\wedge}$）上的作用，来自 $\tilde{SO}_{n}(K)$ 在 $\operatorname{Pf}_{\mathcal{L}}^{\wedge}$ 上的（显然唯一的）作用。
\item $\operatorname{Aut}_{2}O$ 与 $\tilde{SO}_{n}(K)$ 在 $\operatorname{Pf}_{\mathcal{L}}^{\wedge}$ 上的作用，定义了 $\operatorname{Aut}_{2}O\ltimes\tilde{SO}_{n}(K)$ 的一个作用。
\end{enumerate}

\emph{注.} 陈述 (ii) 可按 2.8.2 的精神来理解：$\operatorname{Aut}_{2}O$ 的作用给出 $\pi^{*}\operatorname{Pf}_{\mathcal{L}}$ 上沿 $X$ 的联络，其中 $\pi$ 是态射 $M\to\operatorname{Bun}_{G}$。而 $\operatorname{Aut}_{2}O$ 的作用与 $\tilde{SO}_{n}(K)$ 的作用相容，意思是某个中心扩张 $\tilde{J}^{\mathrm{mer}}(SO_{n})_{\mathcal{L}}$ 在 $\pi^{*}\operatorname{Pf}_{\mathcal{L}}$ 上的作用是水平的。

\emph{证明.} 要定义 $\operatorname{Aut}_{2}O\ltimes\tilde{SO}_{n}(K)$ 在 $\operatorname{Pf}_{\mathcal{L}}^{\wedge}$ 上具有所需性质的作用，我们按 4.3.7 的办法进行。设 $R$ 是 $\mathbb{C}$-代数。取 $M_{2}^{\wedge}$ 的一个 $R$-点 $u$，以及 $\operatorname{Aut}_{2}O\ltimes\tilde{SO}_{n}(K)$ 的一个 $R$-点 $\tilde{g}$。回忆 $SO_{n}$ 是 $SO(W)$ 的缩写。记 $\mathcal{F}$ 与 $\mathcal{F}'$ 为 $X\otimes R$ 上对应 $u$ 与 $gu$ 的 $SO(W)$-torsor，其中 $g$ 是 $\tilde{g}$ 在 $\operatorname{Aut}_{2}O\ltimes SO_{n}(K)$ 中的像。我们须定义一个同构
\begin{equation}\tag{206}
\operatorname{Pf}(W_{\mathcal{F}}\otimes\mathcal{L}_{R})\xrightarrow{\ \sim\ }\operatorname{Pf}(W_{\mathcal{F}'}\otimes\mathcal{L}_{R})
\end{equation}
其中 $\mathcal{L}_{R}$ 是 $\mathcal{L}$ 到 $X\otimes R$ 的拉回。

\SourcePage{163}

置 $V:=\omega_{O}^{1/2}\otimes_{O}K\otimes W$。这是 $\mathbb{C}$ 上的 Tate 空间，带有非退化对称双线性型（见 4.3.3）。由 4.2.14，
\begin{equation}\tag{207}
\operatorname{Pf}(W_{\mathcal{F}}\otimes\mathcal{L}_{R})=\operatorname{Pf}(V\hat\otimes R;L_{+}\hat\otimes R,L^{u}_{-})
\end{equation}
其中 $L_{+}:=\omega_{O}^{1/2}\otimes W\subset V$（故 $L_{+}$ 是 $V$ 中的 Lagrange c-格），而 Lagrange d-格 $L^{u}_{-}\subset V\hat\otimes R$ 定义如下。点 $u\in M_{2}^{\wedge}(R)$ 是四元组 $(\alpha,\mathcal{F},\gamma,f)$，其中 $\alpha$、$\mathcal{F}$、$\gamma$ 的含义同 2.8.4（在我们这个特殊情形 $G=SO(W)$ 中），而 $f$ 是 $R\hat\otimes\omega_{O}$ 的平方根构成的群胚中、$H^{0}(\operatorname{Spec}R\hat\otimes K,\alpha^{*}\mathcal{L}_{R})$ 与 $R\hat\otimes\omega_{O}^{1/2}$ 之间的一个同构。记 $\Gamma_{\alpha}$ 的含义同 2.8.4。则
\[
L^{u}_{-}:=H^{0}((X\otimes R)\setminus\Gamma_{\alpha},W_{\mathcal{F}}\otimes\mathcal{L}_{R})\subset H^{0}(\operatorname{Spec}R\hat\otimes K,\alpha^{*}W_{\mathcal{F}}\otimes\alpha^{*}\mathcal{L}_{R})\xrightarrow[\ \sim\ ]{\varphi}V\hat\otimes R
\]
（同构 $\varphi$ 由 $\gamma$ 与 $f$ 诱导）。

计入 (207)，可见构造 (206) 等价于定义一个同构
\begin{equation}\tag{208}
\operatorname{Pf}(V\hat\otimes R;L_{+}\hat\otimes R,L^{u}_{-})\xrightarrow{\ \sim\ }\operatorname{Pf}(V\hat\otimes R;L_{+}\hat\otimes R,L^{gu}_{-}).
\end{equation}
群归纳概形 $\operatorname{Aut}_{2}O\ltimes SO(W\otimes K)$ 以显然的方式作用在 $V$ 上，且容易看出 $L^{gu}_{-}=gL^{u}_{-}$。由 (166)，(208) 的左端等于 $(M\otimes R)_{L_{-}}$ 的逆，其中 $M$ 是 Clifford 模 $Cl(V)/Cl(V)L_{+}$，$L_{-}:=L^{u}_{-}$。于是只剩构造同构 $(M\otimes R)_{L_{-}}\xrightarrow{\ \sim\ }(M\otimes R)_{gL_{-}}$。我们把它定义为 $\tilde{g}$ 在 $M\otimes R$ 上的作用\footnote{回忆 $g$ 是 $\operatorname{Aut}_{2}O\ltimes\tilde{SO}_{n}(K)=\operatorname{Aut}_{2}O\ltimes\tilde{SO}(W\otimes K)$ 的一个 $R$-点。按 $\tilde{SO}_{n}(K)$ 的定义，它作用在 $M$ 上。群归纳概形 $\operatorname{Aut}_{2}O$ 作用在 $(V,L_{+})$ 上，从而作用在 $M$ 上。}所诱导的同构。\hfill$\Box$

\bdsec{4.4. $\operatorname{Bun}_G$ 上的半形式}

\bdp{4.4.1} 设 $G$ 半单。固定 $\mathfrak{g}$ 上一个 $G$-不变的非退化对称双线性型。置 $n:=\dim\mathfrak{g}$，并把 $SO(\mathfrak{g})$ 写成 $SO_{n}$。伴随表示 $G\to SO(\mathfrak{g})$ 诱导态射 $f:\operatorname{Bun}_{G}\to\operatorname{Bun}_{SO_{n}}$。对 $\mathcal{L}\in\omega^{1/2}(X)$ 置 $\lambda'_{\mathcal{L}}:=f^{*}\operatorname{Pf}_{\mathcal{L}}$，其中 $\operatorname{Pf}_{\mathcal{L}}$ 是 4.3.1 中的线丛；

\SourcePage{164}

故 $\lambda'_{\mathcal{L}}$ 在 $\mathfrak{F}\in\operatorname{Bun}_{G}$ 处的纤维等于 $\operatorname{Pf}(\mathfrak{g}_{\mathfrak{F}}\otimes\mathcal{L})\otimes\operatorname{Pf}(\mathfrak{g}\otimes\mathcal{L})^{\otimes-1}$。同构 (189) 诱导一个同构
\begin{equation}\tag{209}
(\lambda'_{\mathcal{L}})^{\otimes 2}=\omega^{\sharp}_{\operatorname{Bun}_{G}}
\end{equation}
这里 $\omega^{\sharp}_{\operatorname{Bun}_{G}}$ 是正规化典范丛 (146)；按 2.1.1，它在 $\mathfrak{F}\in\operatorname{Bun}_{G}$ 处的纤维等于 $\det R\Gamma(X,\mathfrak{g}_{\mathfrak{F}})\otimes(\det R\Gamma(X,\mathfrak{g}\otimes O_{X}))^{\otimes-1}$。

\bdp{4.4.2} 考虑函子
\begin{equation}\tag{210}
\lambda':\omega^{1/2}(X)\to(\omega^{\sharp})^{1/2}(\operatorname{Bun}_{G}),
\end{equation}
$\mathcal{L}\mapsto\lambda'_{\mathcal{L}}$。由 4.3.10，$\lambda'$ 关于 Picard 函子 $\tilde\ell':\mu_{2}\operatorname{tors}(X)\to\mu_{2}\operatorname{tors}(\operatorname{Bun}_{G})$ 是仿射的；该函子把 $X$ 上的 $\mu_{2}$-torsor $\mathcal{E}$ 送到 $\tilde\ell'_{\mathcal{E}}$，即 $\operatorname{Bun}_{SO_{n}}$ 上的 torsor $\ell^{\operatorname{Spin}}_{\mathcal{E}}$ 到 $\operatorname{Bun}_{G}$ 的拉回。

\bdp{4.4.3} \emph{命题.} $\tilde\ell'=\ell'$，其中 $\ell'$ 是下述两个函子的复合：由 (56) 诱导的函子 $\mu_{2}\operatorname{tors}(X)\to Z\operatorname{tors}(X)$，以及 4.1.1–4.1.4 中构造的函子 $\ell:Z\operatorname{tors}(X)\to\mu_{\infty}\operatorname{tors}(\operatorname{Bun}_{G})$。这里 $Z=\pi_{1}(G)^{\vee}=$ ${}^{L}G$ 的中心（见 4.1.1 的注）。

假定该命题，我们定义典范的 $\ell$-仿射函子
\begin{equation}\tag{211}
\lambda:Z\operatorname{tors}_{\theta}(X)\to\mu_{\infty}\operatorname{tors}_{\theta}(\operatorname{Bun}_{G})
\end{equation}
为 $\mathcal{E}\cdot\mathcal{L}\mapsto\lambda_{\mathcal{E}\cdot\mathcal{L}}:=\ell_{\mathcal{E}}\cdot\lambda'_{\mathcal{L}}$，其中 $\mathcal{E}\in Z\operatorname{tors}(X)$，$\mathcal{L}\in\omega^{1/2}(X)$。（注意：正规化问题！！！？？？）
% [Tutor Note: 原书 p.164 此括注「Attention: normalization problem!!!???」是作者本人的自注，连用感叹号与问号，底本如此，保留原样。]

为证明命题 4.4.3，注意 $\tilde\ell'$ 是对应于 $G$ 被 $\mu_{2}$ 的扩张的函子 (152)，该扩张由 $SO(\mathfrak{g})$ 的旋量扩张诱导。因此 $\tilde\ell'$ 是 $\ell:Z\operatorname{tors}(X)\to\mu_{\infty}\operatorname{tors}(\operatorname{Bun}_{G})$ 与函子 $\mu_{2}\operatorname{tors}(X)\to Z\operatorname{tors}(X)$ 的复合，后者由态射 $\mu_{2}\to Z=\pi_{1}(G)^{\vee}$ 诱导，而该态射对偶于 $\pi_{1}(G)\to\pi_{1}(SO(\mathfrak{g}))=\mathbb{Z}/2\mathbb{Z}$。于是只需证明下面的事实。

\SourcePage{165}

\bdp{4.4.4} \emph{引理.} 态射 $\pi_{1}(G)\to\pi_{1}(SO(\mathfrak{g}))=\mathbb{Z}/2\mathbb{Z}$ 对偶于群 ${}^{L}G$ 的态射 (56)。

\emph{证明.} 有典范同构 $f:P/P_{G}\xrightarrow{\ \sim\ }\operatorname{Hom}(\pi_{1}(G)(1),\mu_{\infty})$，其中 $P_{G}$ 是 $G$ 的权群，$P$ 是其泛覆叠 $\tilde{G}$ 的权群；权 $\lambda\in P$ 是嘉当子群 $\tilde{H}\subset\tilde{G}$ 的特征标，而 $f(\lambda)$ 是它在 $\pi_{1}(G)(1)\subset\tilde{H}$ 上的限制。设 $M$ 是 $so(\mathfrak{g})$ 的旋量表示。则 $\tilde{G}$ 作用在 $M$ 上，且 $\pi_{1}(G)(1)\subset\tilde{G}$ 按某个特征标 $\chi\in\operatorname{Hom}(\pi_{1}(G)(1),\mu_{\infty})$ 作用。按 (56) 的定义（另见 3.4.1 中 $\lambda^{\#}$ 的定义），引理只是说 $\chi=f(\rho)$，其中 $\rho\in P$ 是基本权之和。

设 $\mathfrak{b}\subset\mathfrak{g}$ 是 Borel 子代数。选取 $\mathfrak{b}$-不变的旗 $0\subset V_{1}\subset\ldots\subset V_{n}=\mathfrak{g}$，使 $\dim V_{k}=k$、$V_{k}^{\perp}=V_{n-k}$，且 $\mathfrak{b}$ 是其中某个 $V_{k}$。设 $\mathfrak{b}'$ 是该旗在 $so(\mathfrak{g})$ 中的稳定化子。这是 $so(\mathfrak{g})$ 中含 $\mathfrak{b}$ 的 Borel 子代数。设 $m\in M$ 是关于 $\mathfrak{b}'$ 的最高向量。则 $\mathbb{C}m$ 是 $\mathfrak{b}$-不变的，且相应的 $\mathfrak{b}$ 的特征标等于正根之和的一半，即 $\rho$。故 $\chi=f(\rho)$。\hfill$\Box$

\emph{注.} 按 Kostant（参看 [Ko61] 中引理 5.9 的证明），$\mathfrak{g}$-模 $M$ 同构于 $2^{[r/2]}$ 个最高权为 $\rho$ 的不可约 $\mathfrak{g}$-模之和（其中 $r$ 是 $\mathfrak{g}$ 的秩）。

\bdp{4.4.5} 我们对 (211) 的构造略依赖 $\mathfrak{g}$ 上标量积的选取（见 4.4.1）。由于 $\mathfrak{g}$ 上有若干「典范」标量积，读者或许更喜欢 (211) 的如下版本。

为简化记号，设 $G$ 单。则 $\mathfrak{g}$ 上不变对称双线性型的空间是一维的。记之为 $\beta$。取 $\beta$ 的一个平方根，即一维向量空间 $\beta^{1/2}$，连同一个同构 $\beta^{1/2}\otimes\beta^{1/2}\xrightarrow{\ \sim\ }\beta$。于是 $\mathfrak{g}\otimes\beta^{1/2}$ 带有典范双线性型。考虑表示 $G\to SO(\mathfrak{g}\otimes\beta^{1/2})$，然后按 4.4.1–4.4.3 进行（例如，此时 $\lambda'_{\mathcal{L}}$ 在 $\mathcal{F}\in\operatorname{Bun}_{G}$ 处的纤维等于 $\operatorname{Pf}(\mathfrak{g}_{\mathcal{F}}\otimes\mathcal{L}\otimes\beta^{1/2})\otimes\operatorname{Pf}(\mathfrak{g}\otimes\mathcal{L}\otimes\beta^{1/2})^{\otimes-1}$）。这样得到的函子 (211)

\SourcePage{166}

略依赖 $\beta^{1/2}$ 的选取。更确切地说，$-1\in\operatorname{Aut}\beta^{1/2}$ 作用在 $\lambda'_{\mathcal{L}}$ 上，从而作用在 $\lambda_{\mathcal{M}}$（$\mathcal{M}\in Z\operatorname{tors}_{\theta}(X)$）上，其方式为乘以 $(-1)^{p}$，其中 $p:\operatorname{Bun}_{G}\to\mathbb{Z}/2\mathbb{Z}$ 是复合
\[
\operatorname{Bun}_{G}\to\pi_{0}(\operatorname{Bun}_{G})=\pi_{1}(G)\to\pi_{1}(SO(\mathfrak{g}))=\mathbb{Z}/2\mathbb{Z}.
\]
我们是否该把 $\lambda_{\mathcal{M}}$ 看作一个 SUPER-层？？？！
% [Tutor Note: 原书 p.166「Do we want to consider $\lambda_{\mathcal{M}}$ as a SUPER-sheaf???!」是作者行文中的自问，连用问号与感叹号，底本如此，保留原样。]

\bdp{4.4.6} 我们已把 $\mathcal{L}\in Z\operatorname{tors}_{\theta}(X)$ 对应到 $\operatorname{Bun}_{G}$ 上的线丛 $\lambda_{\mathcal{L}}$（见 4.4.1–4.4.3）。对 $x\in X$，记 $\lambda_{\mathcal{L},\underline{x}}$ 为 $\lambda_{\mathcal{L}}$ 到 $\operatorname{Bun}_{G,\underline{x}}$ 的拉回。在 4.4.7–4.4.10 中我们将定义 $G(K_{x})$ 的中心扩张 $\tilde{G}(K_{x})_{\mathcal{L}}$，它作用在 $\lambda_{\mathcal{L},\underline{x}}$ 上。在 4.4.11–4.4.13 中我们考虑 $\tilde{G}(K_{x})_{\mathcal{L}}$ 的李代数。

\bdp{4.4.7} 设 $O$、$K$ 与 $\omega_{O}$ 的含义同 4.3.3。固定 $\omega_{O}$ 的一个平方根 $\mathcal{L}$。然后我们构造群归纳概形的中心扩张
\begin{equation}\tag{212}
0\to\mathbb{G}_{m}\to\tilde{G}(K)_{\mathcal{L}}\to G(K)\to 0
\end{equation}
如下。$\mathcal{L}$ 定义了中心扩张 (196)。固定 $\mathfrak{g}$ 上一个非退化不变对称双线性型\footnote{不固定 $\mathfrak{g}$ 上的型，读者也可按 4.4.5 处理。}，并把 $SO(\mathfrak{g})$ 写成 $SO_{n}$，$n:=\dim\mathfrak{g}$。我们把 (212) 定义为 $G(K)$ 的中心扩张，它与由 (196) 经伴随表示 $G\to SO(\mathfrak{g})=SO_{n}$ 诱导的扩张相反。扩张 (212) 在 $G(O)$ 上分裂。

\emph{注.} 在 $G=SO_{r}$ 的情形，我们的记号有歧义：$\tilde{G}(K)\ne\tilde{SO}_{r}(K)$。但愿这一歧义无害。

\bdp{4.4.8} 设 $\mathcal{L}\in\omega^{1/2}(X)$，$x\in X$。按 4.4.7，$\mathcal{L}$ 到 $\operatorname{Spec}O_{x}$ 的限制定义了 $G(K_{x})$ 的中心扩张，记之为 $\tilde{G}(K_{x})_{\mathcal{L}}$。记 $\lambda'_{\mathcal{L},\underline{x}}$ 为线丛 $\lambda'_{\mathcal{L}}$（来自 4.4.1）到 $\operatorname{Bun}_{G,\underline{x}}$ 的拉回。由 4.3.7，$G(K_{x})$ 在 $\operatorname{Bun}_{G,\underline{x}}$ 上的作用提升为 $\tilde{G}(K_{x})_{\mathcal{L}}$ 在 $\lambda'_{\mathcal{L}}$ 上的典范作用。
% [Tutor Note: 原书 p.166 末三处均印作 $\lambda'_{\mathcal{L}}$（无逗号与下划线 x 的下标），而同段前文定义的是 $\lambda'_{\mathcal{L},\underline{x}}$；底本如此，保留原样。]

子群 $\mathbb{G}_{m}\subset\tilde{G}(K_{x})_{\mathcal{L}}$ 以显然的方式作用在 $\lambda'_{\mathcal{L}}$ 上（见 4.4.7 中 $\tilde{G}(K_{x})_{\mathcal{L}}$ 的定义，以及 4.3.7 的最后

\SourcePage{167}

一句）。$G(O_{x})\subset\tilde{G}(K_{x})_{\mathcal{L}}$ 在 $\lambda'_{\mathcal{L},\underline{x}}$ 上的作用是显然的那个。

\bdp{4.4.9} 在 4.4.7 中我们定义了一个函子
\begin{equation}\tag{213}
\omega^{1/2}(O)\to\{G(K)\text{ 被 }\mathbb{G}_{m}\text{ 的中心扩张}\}
\end{equation}
其中 $\omega^{1/2}(O)$ 是 $\omega_{O}$ 的平方根构成的群胚。(213) 的左端是 3.4.4 意义下的 $\mu_{2}$-范畴。(213) 的右端是 $Z$-范畴，$Z:=\pi_{1}(G)^{\vee}=\operatorname{Hom}(\pi_{1}(G),\mathbb{G}_{m})$。事实上，边缘态射\footnote{先验地，(214) 是抽象群的态射；但按 4.1.7 的注，它实际上是群归纳概形的态射。另见 4.5.4。}
\begin{equation}\tag{214}
G(K)\to H^{1}(K,\pi_{1}^{\mathrm{et}}(G))=\pi_{1}(G)=Z^{\vee}
\end{equation}
诱导一个态射\footnote{事实上是同构（见 4.5.4）。}
\begin{equation}\tag{215}
Z\to\operatorname{Hom}(G(K),\mathbb{G}_{m}),
\end{equation}
即 (213) 右端上的一个 $Z$-结构。用 (56) 定义的态射 $\mu_{2}\to Z$，我们把 (213) 的右端看作 $\mu_{2}$-范畴。于是 (213) 是 $\mu_{2}$-函子（用 4.3.4、4.3.4 的注 (ii) 以及 4.4.4）。故由 3.4.4，函子 (213) 给出一个 $Z$-函子
\begin{equation}\tag{216}
Z\operatorname{tors}_{\theta}(O)\to\{G(K)\text{ 被 }\mathbb{G}_{m}\text{ 的中心扩张}\}.
\end{equation}
由 (213) 给出的、对应于 $\mathcal{L}\in Z\operatorname{tors}_{\theta}(O)$ 的 $G(K)$ 的中心扩张将记作 $\tilde{G}(K)_{\mathcal{L}}$。扩张
\begin{equation}\tag{217}
0\to\mathbb{G}_{m}\to\tilde{G}(K)_{\mathcal{L}}\to G(K)\to 0
\end{equation}
在 $G(O)$ 上分裂。

\emph{注记.}
\begin{enumerate}[label=(\roman*),leftmargin=2.2em]
\item 按 3.4.7 (i)，(213) 右端上的 $Z$-结构给出一个 Picard 函子
\begin{equation}\tag{218}
Z\operatorname{tors}(O)=Z\operatorname{tors}\to\{G(K)\text{ 被 }\mathbb{G}_{m}\text{ 的中心扩张}\}.
\end{equation}
\end{enumerate}

\SourcePage{168}

显式地，(218) 是如下两个函子的复合。其一是典范等价
\begin{equation}\tag{219}
\{\mathbb{G}_{m}\text{ 对 }Z^{\vee}\text{ 的平凡扩张}\}=Z\operatorname{tors},\qquad
\text{扩张}\mapsto\text{其分裂构成的 }Z\text{-torsor}
\end{equation}
其二是从 (219) 左端到 (218) 右端的函子，它由 (214) 诱导。换言之，(218) 就是 4.1.8 中的函子 $\mathcal{E}\mapsto\tilde{G}(K)_{\mathcal{E}}$。

(ii) 由 3.4.7 (iv)，函子 (216) 关于 Picard 函子 (218) 是仿射的。

\bdp{4.4.10} 设 $\mathcal{L}\in Z\operatorname{tors}_{\theta}(X)$。按 4.4.9，$\mathcal{L}$ 在 $Z\operatorname{tors}_{\theta}(O_{x})$ 中的像定义了 $G(K_{x})$ 的中心扩张，记之为 $\tilde{G}(K_{x})_{\mathcal{L}}$。记 $\lambda_{\mathcal{L},\underline{x}}$ 为 $\lambda_{\mathcal{L}}$ 到 $\operatorname{Bun}_{G,\underline{x}}$ 的拉回。$G(K_{x})$ 在 $\operatorname{Bun}_{G,\underline{x}}$ 上的作用提升为 $\tilde{G}(K_{x})_{\mathcal{L}}$ 在 $\lambda_{\mathcal{L},\underline{x}}$ 上的典范作用（用 4.3.7–4.3.9、4.1.8 以及 4.4.9 的注记）。$G(O_{x})\times\mathbb{G}_{m}\subset\tilde{G}(K_{x})_{\mathcal{L}}$ 以显然的方式作用在 $\lambda_{\mathcal{L},\underline{x}}$ 上。

\bdp{4.4.11} \emph{命题.} 对应于 (217) 的李代数扩张是 2.5.1 中的扩张
\[
0\to\mathbb{C}\to\widetilde{\mathfrak{g}}\otimes K\to\mathfrak{g}\otimes K\to 0.
\]

\emph{证明.} 对应于 (217) 的李代数扩张不依赖 $\mathcal{L}\in Z\operatorname{tors}_{\theta}(O)$，故可用 (212) 代替 (217)，最终用 (194)。于是只需用 Kac–Peterson–Frenkel 定理，它说李代数扩张
\begin{equation}\tag{220}
0\to\mathbb{C}\to\tilde{o}_{n}(K)\to o_{n}(K)\to 0
\end{equation}
对应于 (194)，由上闭链 $(u,v)\mapsto\frac{1}{2}\operatorname{Res}\operatorname{Tr}(du,v)$ 定义，$u,v\in o_{n}(K)$。事实上，要用 [KP] 或 [Fr81] 中的命题 I.3.11，须用到 $\tilde{o}_{n}(K)$ 的如下刻画（它不涉及群 $\tilde{O}_{n}(K)$）：设 $V$ 的含义同 4.3.3，$M$ 是 $Cl(V)$ 上的不可约离散模，则有

\SourcePage{169}

$\tilde{o}_{n}(K)$ 在 $M$ 中的表示，它与 $\tilde{o}_{n}(K)$ 在 $Cl(V)$ 上的作用相容，且使 $1\in\mathbb{C}\subset\tilde{o}_{n}(K)$ 恒等地作用在 $M$ 上。\hfill$\Box$

\bdp{4.4.12} 设 $\lambda_{\mathcal{L}}$ 与 $\lambda_{\mathcal{L},\underline{x}}$ 的含义同 4.4.10。按 4.4.10 与 4.4.11，$\mathfrak{g}\otimes K_{x}$ 在 $\operatorname{Bun}_{G,\underline{x}}$ 上的作用提升为 $\widetilde{\mathfrak{g}\otimes K_{x}}$ 在 $\lambda_{\mathcal{L},\underline{x}}$ 上的典范作用；它在 $\mathbb{C}\times(\mathfrak{g}\otimes O_{x})\subset\widetilde{\mathfrak{g}\otimes K_{x}}$ 上的限制是显然的那个。特别地，$1\in\mathbb{C}\subset\widetilde{\mathfrak{g}\otimes K_{x}}$ 作用为乘以 1。

$\lambda_{\mathcal{L}}$ 带有一个同构 $\lambda_{\mathcal{L}}^{\otimes 2n}\xrightarrow{\ \sim\ }(\omega^{\sharp}_{\operatorname{Bun}_{G}})^{\otimes n}$（对某个 $n\ne 0$），故作用在 $\lambda_{\mathcal{L}}$ 上的微分算子层是 $D'$。因此由 1.2.5，$\widetilde{\mathfrak{g}\otimes K_{x}}$ 在 $\lambda_{\mathcal{L},\underline{x}}$ 上的作用诱导一个典范态射
\[
h_{x}:\mathfrak{z}_{x}\to\Gamma(\operatorname{Bun}_{G},D').
\]
显然 $h_{x}$ 不依赖 $\mathcal{L}\in Z\operatorname{tors}_{\theta}(X)$。

\bdp{4.4.13} 本小节我们证明 4.4.12 中的 $h_{x}$ 与 2.5.4 中的 $h_{x}$ 一致。读者可跳过此证明，干脆忘掉 $h_{x}$ 的旧定义（当初引入它，只是为了在 §2 中回避 $\omega_{\operatorname{Bun}_{G}}$ 的平方根的讨论）。

要证 $h_{x}$ 的两个定义等价，只需说明：若 $\mathcal{L}$ 是 $\omega_{X}$ 的平方根，则同构 $\lambda_{\mathcal{L}}^{\otimes 2}\xrightarrow{\ \sim\ }\omega^{\sharp}_{\operatorname{Bun}_{G}}$ 在它们到 $\operatorname{Bun}_{G,\underline{x}}$ 的拉回之间诱导一个 $\mathfrak{g}\otimes K_{x}$-等变同构。这可以直接证明，但事实上不可能有别的情形。事实上，$\widetilde{\mathfrak{g}\otimes K_{x}}$-等变性的障碍是一个上闭链 $\widetilde{\mathfrak{g}\otimes K_{x}}\to H^{0}(\operatorname{Bun}_{G,\underline{x}},\mathcal{O})$。因 $\operatorname{Hom}(\mathfrak{g}\otimes K_{x},\mathbb{C})=0$，只需说明 $\operatorname{Bun}_{G,\underline{x}}$ 上每个正则函数 $f$ 局部常值。按 2.3.1，$\operatorname{Bun}_{G,\underline{x}}$ 是 $\operatorname{Bun}_{G,nx}$（$n\in\mathbb{N}$）的逆向极限。显然 $f$ 来自某个 $\operatorname{Bun}_{G,nx}$ 上的正则函数。故只需证明下述引理。

\emph{引理.} $\operatorname{Bun}_{G,nx}$ 上每个正则函数局部常值。

\emph{证明.} 取 $y\in X\setminus\{x\}$，考虑概形 $M$，它参数化 $X$ 上在 $nx$ 与 $y$ 的形式邻域上都平凡化的 $G$-丛（这里除子 $nx$ 看作子概形）。$G(K_{y})$ 作用在 $M$ 上，而正则

\SourcePage{170}

函数 $f$ 是 $H^{0}(M,\mathcal{O}_{M})$ 的 $G(O_{y})$-不变元。显然 $H^{0}(M,\mathcal{O}_{M})$ 是可积离散 $\mathfrak{g}\otimes K_{y}$-模。众所周知且极易证明：这样的模中 $(\mathfrak{g}\otimes O_{y})$-不变元必是 $(\mathfrak{g}\otimes K_{y})$-不变元。故 $f$ 是 $(\mathfrak{g}\otimes K_{y})$-不变的。因 $\mathfrak{g}\otimes K_{y}$ 在 $M$ 上的作用是（形式地）传递的，$f$ 局部常值。\hfill$\Box$

\emph{注.} 上述引理是众所周知的。标准的证法是把 $\operatorname{Bun}_{G,nx}$ 表示成 $\Gamma\backslash G(K_{y})/G(O_{y})$（对某个 $\Gamma\subset G(K_{y})$；$n=0$ 的情形见 [La-So]），然后利用 $G(K_{y})/G(O_{y})$ 上正则函数局部常值这一事实。

\bdp{4.4.14} 现在我们来完成水平性定理 2.7.3 的证明（证明的开头见 2.8.3–2.8.5）。

设 $M$ 是 $X$ 上的概形，它在 $x\in X$ 处的纤维是 $\operatorname{Bun}_{G,\underline{x}}$。固定 $\mathcal{L}\in\omega^{1/2}(X)$ 与 $\mathcal{L}^{\mathrm{loc}}\in\omega^{1/2}(O)$（即 $\mathcal{L}$ 是 $\omega_{X}$ 的平方根，$\mathcal{L}^{\mathrm{loc}}$ 是 $\omega_{O}$ 的平方根）。则有 4.3.16 中定义的概形 $X_{2}^{\wedge}$。记 $M_{2}^{\wedge}$ 为 $M$ 与 $X_{2}^{\wedge}$ 在 $X$ 上的纤维积。半直积 $\operatorname{Aut}_{2}O\ltimes G(K)$ 作用在 $M_{2}^{\wedge}$ 上（参看 4.3.16）。

取其中心扩张 $\operatorname{Aut}_{2}O\ltimes\widetilde{G(K)}$，其中 $\widetilde{G(K)}$ 是对应于 $\mathcal{L}^{\mathrm{loc}}$ 的中心扩张 (212)，而 $\operatorname{Aut}_{2}O=\operatorname{Aut}(O,\mathcal{L}^{\mathrm{loc}})$ 按结构的搬运作用在 $\widetilde{G(K)}=\widetilde{G(K)}_{\mathcal{L}^{\mathrm{loc}}}$ 上。记 $\lambda_{\mathcal{L}}^{\wedge}$ 为 Pfaffian 线丛 $\lambda'_{\mathcal{L}}$（来自 4.4.1）到 $M_{2}^{\wedge}$ 的拉回。因 $\operatorname{Aut}_{2}O$ 作用在 $M_{2}^{\wedge}$ 上，而 $M_{2}^{\wedge}$ 看作 $\operatorname{Bun}_{G}$ 上的概形，我们得到 $\operatorname{Aut}_{2}O$ 在 $\lambda_{\mathcal{L}}^{\wedge}$ 上的作用。另一方面，$\widetilde{G(K)}$ 作用在 $\lambda_{\mathcal{L},\hat x}^{\wedge}$ 上，后者定义为 $\lambda_{\mathcal{L}}^{\wedge}$ 在 $M_{2}^{\wedge}$ 于 $\hat x\in X_{2}^{\wedge}$ 处纤维上的限制。事实上，该纤维等于 $\operatorname{Bun}_{G,\underline{x}}$，其中 $x$ 是 $\hat x$ 在 $X$ 中的像；由 4.4.8，中心扩张 $\tilde{G}(K_{x})_{\mathcal{L}}$ 作用在 $\lambda'_{\mathcal{L},\underline{x}}=\lambda_{\mathcal{L},\hat x}^{\wedge}$ 上。该扩张只依赖 $\mathcal{L}_{x}$，即 $\mathcal{L}$ 到 $\operatorname{Spec}O_{x}$ 的拉回。因 $\hat x$ 给出同构 $(O_{x},\mathcal{L}_{x})\xrightarrow{\ \sim\ }(O,\mathcal{L}^{\mathrm{loc}})$，我们得到同构 $\tilde{G}(K_{x})_{\mathcal{L}}\xrightarrow{\ \sim\ }\widetilde{G(K)}$，从而得到 $\widetilde{G(K)}$ 在 $\lambda_{\mathcal{L},\hat x}^{\wedge}$ 上的作用。如 2.8.5 所述，为完成 2.7.3 的证明，只需说明

\SourcePage{171}

(i) 对应于各种 $\hat x\in X_{2}^{\wedge}$ 的 $\widetilde{G(K)}$ 在 $\lambda_{\mathcal{L},\hat x}^{\wedge}$ 上的作用，来自 $\widetilde{G(K)}$ 在 $\lambda_{\mathcal{L}}^{\wedge}$ 上的（显然唯一的）作用。

(ii) 该作用与 $\operatorname{Aut}_{2}O$ 的作用相容（即实际上我们有 $\operatorname{Aut}_{2}O\ltimes\widetilde{G(K)}$ 在 $\lambda_{\mathcal{L}}^{\wedge}$ 上的作用）。

这立即由 4.3.17 推出。

\bdp{4.4.15} 本小节与下一小节中，我们陈述并证明 4.4.14 的陈述 i) 与 ii) 的一个推广，它将在本书主要结果（定理 5.4.5）的证明中用到。推广是显然的（把 $\omega^{1/2}(X)$ 换成 $Z\operatorname{tors}_{\theta}(X)$，等等），读者目前完全可以跳过这两小节。

固定 $\mathcal{L}\in Z\operatorname{tors}_{\theta}(X)$ 与 $\mathcal{L}^{\mathrm{loc}}\in Z\operatorname{tors}_{\theta}(O)$。记 $X_{Z}^{\wedge}$ 为 $X^{\wedge}$ 的平展 $Z$-覆叠，使得 $X_{Z}^{\wedge}(R)$ 中对应于态射 $\alpha:\operatorname{Spec}(R\hat\otimes O)\to X$ 的点的原像，是群胚 $Z\operatorname{tors}_{\theta}(R\hat\otimes O)$ 中同构 $\mathcal{L}_{R}^{\mathrm{loc}}\xrightarrow{\ \sim\ }\alpha^{*}\mathcal{L}$ 的集合\footnote{这里为方便起见，采用 3.4.5 中 $Z\operatorname{tors}_{\theta}$ 的定义。}，其中 $\mathcal{L}_{R}^{\mathrm{loc}}$ 是 $\mathcal{L}^{\mathrm{loc}}$ 到 $\operatorname{Spec}R\hat\otimes O$ 的拉回。4.6.6 中的群归纳概形 $\operatorname{Aut}_{Z}O=\operatorname{Aut}(O,\mathcal{L}^{\mathrm{loc}})$ 按结构的搬运作用在 $X_{Z}^{\wedge}$ 上。记 $M_{Z}^{\wedge}$ 为 $M$ 与 $X_{Z}^{\wedge}$ 在 $X$ 上的纤维积。记 $\lambda_{\mathcal{L}}^{\wedge}$ 为 4.4.3 中定义的线丛 $\lambda_{\mathcal{L}}$ 到 $M_{Z}^{\wedge}$ 的拉回。半直积 $\operatorname{Aut}_{Z}O\ltimes G(K)$ 作用在 $M_{Z}^{\wedge}$ 上。取其中心扩张 $\operatorname{Aut}_{Z}O\ltimes\widetilde{G(K)}$，其中 $\widetilde{G(K)}$ 是对应于 $\mathcal{L}^{\mathrm{loc}}$ 的中心扩张 (217)，而 $\operatorname{Aut}_{Z}O=\operatorname{Aut}(O,\mathcal{L}^{\mathrm{loc}})$ 按结构的搬运作用在 $\widetilde{G(K)}=\widetilde{G(K)}_{\mathcal{L}^{\mathrm{loc}}}$ 上。我们把 $\operatorname{Aut}_{Z}O\ltimes G(K)$ 在 $M_{Z}^{\wedge}$ 上的作用提升为 $\operatorname{Aut}_{Z}O\ltimes\widetilde{G(K)}$ 在 $\lambda_{\mathcal{L}}^{\wedge}$ 上的作用。

正如 4.4.14，可定义 $\operatorname{Aut}_{Z}O$ 在 $\lambda_{\mathcal{L}}^{\wedge}$ 上的作用，以及 $\widetilde{G(K)}$ 在 $\lambda_{\mathcal{L},\hat x}^{\wedge}:=$ $\lambda_{\mathcal{L}}^{\wedge}$ 在 $M_{Z}^{\wedge}$ 于 $\hat x\in\hat X_{Z}$ 处纤维上的限制上的作用。
% [Tutor Note: 原书 p.171 同段前文印作 $X_{Z}^{\wedge}$，此处末句印作 $\hat X_{Z}$；底本如此，保留原样。]

\bdp{4.4.16} \emph{命题.}

(i) 对应于各种 $\hat x\in X_{Z}^{\wedge}$ 的 $\widetilde{G(K)}$ 在 $\lambda_{\mathcal{L},\hat x}^{\wedge}$ 上的作用，来自 $\widetilde{G(K)}$ 在 $\lambda_{\mathcal{L}}^{\wedge}$ 上的（显然唯一的）作用。
